% By Paul E. West and Sean T. Hayes
\documentclass[14pt,aspectratio=1610]{beamer}
% \documentclass[12pt,aspectratio=1610,handout]{beamer}
\usepackage{csu-theme}
\usepackage{booktabs}
\usepackage{forest}
\usetikzlibrary{positioning,backgrounds}

% for aligning items in an itemize (eqmakebox)
\usepackage{eqparbox}

% Add space between paragraphs
\setlength{\parskip}{0.5\baselineskip}

\title{\Huge{}{\ttfamily\bfseries C++} Development\\in Linux}
\subtitle{CSCI 315: \textit{Data Structure Analysis}}
\author{Sean T. Hayes, Ph.D.}
\institute{Charleston Southern University}
%% Comment out date to use default (today)
%\date{January 16, 2025}
\subject{Data Structures}
\keywords{Linux, Unix, C++, Compilation}

\graphicspath{{../imgs/compilation}}

% Change font-family and color for description lists to type-writer
\setbeamerfont{description item}{family=\ttfamily}
\setbeamercolor{description item}{fg=codeVariable, bg=codeBackground}

\begin{document}

\begin{frame}
  \titlepage
\end{frame}

\section{Introduction}

\begin{frame}{Learning Outcomes}
  Linux Philosophy
  \begin{itemize}
    \item Describe the core principles of the Unix/Linux philosophy.
    \item Identify and navigate the Linux file system hierarchy.
    \item Explain the purpose of regular files, directories, and device files in Linux.
  \end{itemize}
\end{frame}

\begin{frame}{Learning Outcomes}
  C++ Software Development
  \begin{itemize}
      \item Explain the purpose and role of a compiler.
      \item Describe the stages of the compilation process in C/C++.
      \item Interpret and use basic g++ command-line options.
      \item Differentiate between source files, object files, and executables for modularized code.
      \item Understand why C++ separates headers (with guards) for interfaces from source files for implementations.
      \item Use make and basic Makefile rules to automate multi-file C++ compilation tasks.
      \item Diagnose common compiler and linker errors.
  \end{itemize}
\end{frame}

\section{Linux Philosophy}

\begin{frame}{The Linux File Structure}
\begin{itemize}
  \item<+-> In Linux, \href{https://en.wikipedia.org/wiki/Everything_is_a_file}{everything is a file}! Well, almost.
  \item<+-> Drives, ports, and devices (e.g., printers) are file descriptors.
\end{itemize}
\end{frame}

\begin{frame}{The Linux File Structure uses Inodes}

\begin{itemize}
  \item<+-> What an \href{https://en.wikipedia.org/wiki/Inode}{\emph{inode}} is:
  \begin{itemize}
    \item index node
    \item administrative metadata
    \item a special block of data in the file system
  \end{itemize}

  \item<+-> 
    What it contains:
    \begin{itemize}
      \item owner
      \item group
      \item permissions
      \item timestamps (creation/modification date)
      \item file size
      \item disk locations
    \end{itemize}
\end{itemize}
\end{frame}

\begin{frame}{Directories}
\begin{itemize}[<+->]
\item Directories (a.k.a.\ folders) are also files.
\item A directory is a file that holds the \emph{inode} (index node) numbers and names of other files.
\item Each directory entry is a link to a file's inode; remove the filename and you remove the link.
\item You can see the inode number for a file by using \lstinline[language=Bash]{ls -i}.
\item If the last link to a file is deleted, the inode and referenced data are marked as free (soft delete).
\item This allows deletion when there are multiple hard links to the same file to be managed correctly.
\end{itemize}
\end{frame}

\begin{frame}{Directories continued}
Files are arranged in directories, which may contain sub-directories.

\begin{tabularx}{1\textwidth}{l X}
\toprule{}
  Directory & Description\\
\midrule{}
  \texttt{/} & Contains all of the system's files in directories.\\
  \texttt{/home} &  Contains a subdirectory for each user's home.\\
  \texttt{/bin}  &  System programs (``binaries'')\\
  \texttt{/etc}  &  System configuration files\\
  \texttt{/lib}  &  System libraries\\
  \texttt{/dev}  &  Physical devices and device interfaces\\
\bottomrule{}
\end{tabularx}

\end{frame}

\begin{frame}{Example Directories Tree}
\centerline{%
  \begin{forest}
    for tree={font=\ttfamily\large, edge=very thick}
    [/
      [bin]
      [etc]
      [dev]
      [home
        [olivia
          [Documents]
          [Downloads]
          [School]
        ]
        [emma]
      ]
      [lib]
    ]
  \end{forest}
}
\end{frame}

\begin{frame}{Files and Devices}
\begin{itemize}
  \item Even hardware devices are very often represented (mapped) by files.
  \item You can mount a CD-ROM drive as a file:
    \begin{itemize}
      \item \lstinline[language=Bash, morekeywords=mount]{mount -t iso9660 /dev/sr0 /mnt/cdrom}
      \item \lstinline[language=Bash]{cd /mnt/cdrom} \# navigate to the mounded drive
    \end{itemize}\vspace{1em}
    \item You can mount a USB drive as a file:
    \begin{itemize}
      \item \lstinline[language=Bash, morekeywords=mkdir]{mkdir /media/usb-drive} \# Make a folder ot be the mount point.
      \item \lstinline[language=Bash, morekeywords=fdisk]{fdisk -l} \# look at the list of available drives
      \item \lstinline[language=Bash, morekeywords=mount]{mount /dev/sdc1 /media/usb-drive/}
      \item \lstinline[language=Bash]{cd /media/usb-drive}
    \end{itemize}
\end{itemize}
\end{frame}

\begin{frame}{\texttt{/dev/console}}
\begin{itemize}
  \item This device represents the system console.
  \item Error messages and diagnostics are often sent to this device.
  \item On Linux, it's usually the ``active'' virtual console.
\end{itemize}
\end{frame}

\begin{frame}{\texttt{/dev/tty}}
\begin{itemize}
  \item The special file \texttt{/dev/tty} is an alias for the controlling terminal of a process.
  \begin{itemize}
    \item keyboard
    \item screen
    \item window
  \end{itemize}
  \item \texttt{/dev/tty} allows a program to write directly to the user, without regard to which pseudo-terminal or hardware terminal the user is using.
\end{itemize}
\end{frame}

\begin{frame}{/dev/null}
\begin{itemize}
\item This is the null device.
\item All output written to this device is discarded.
\item Unwanted output (a.k.a., a student's email complaint/rant) is often redirected to \texttt{/dev/null}.
\item \lstinline[language=bash, deletekeywords=do]{echo I do not want to see this. > /dev/null}
\item \lstinline[language=bash, morekeywords=cp]{cp /dev/null empty_file}
\end{itemize}
\end{frame}

% \section{tmux}

% \begin{frame}{tmux: Terminal Multiplexer}
% \begin{itemize}
% \item This allows us to attach multiple users to the same terminal or to background processes.
% \item For this class, we may use tmux to program as a group or to help everyone follow along.
% \item How to connect:
% \begin{enumerate}
% \item First, everyone must be on the CSU \textit{wired} network.
% \item I will post the IP address in class.
% \item SSH over to that IP address. (You may use PuTTY if you are in Windows.)
% \item Log in is student/student
% \item At the command line type: \\
% \lstinline{tmux att}
% \item To detach from the session, type \lstinline{ctrl + b} and then \lstinline{d}.
% \end{enumerate}
% \end{itemize}
% \end{frame}

\section{Compilation}

\begin{frame}{Introduction}
As a programmer or system administrator, you should know how to program under Linux.

We are going to learn how to use\textellipsis{}\vspace{-0.8em}%
\begin{description}[\texttt{gdb} \& \texttt{ddd~}]
\item[\colorbox{codeBackground}{\texttt{g++}}] to compile a C++ program under Linux.
\item[\colorbox{codeBackground}{\texttt{gdb \& \texttt{ddd}}}] to debug.
\item[\colorbox{codeBackground}{\texttt{make}}] to automate compilation.
\item[\colorbox{codeBackground}{\texttt{valgrind}}] to perform memory analysis.
\item[\colorbox{codeBackground}{\texttt{gnuplot}}] to create performance graphs.
\item[\colorbox{codeBackground}{\texttt{perf}}] to identify performance issues (time permitting).
\end{description}
\end{frame}

\begin{frame}[fragile]{Compiling C++ Programs}
We start with the simple case of a single source-code file.

Create a \texttt{.cpp} file similar to the one listed here.

\begin{lstlisting}
#include <iostream>

int main()
{
    std::cout << "Hello, World!\n";

    return 0;
}
\end{lstlisting}
\end{frame}

\begin{frame}{Compile}
\begin{itemize}
\item Compile using\\
  \lstinline[language={}, basicstyle=\ttfamily\color{codeVariable}]{g++ greeting.cpp}
\item What file is generated?
\item Name your compiled executable by using\\
  \lstinline[language={}, basicstyle=\ttfamily\color{codeVariable}]{g++ greeting.cpp -o welcome}
\item Run the generated executable file.\\
\lstinline[language={}, basicstyle=\ttfamily\color{codeVariable}]{./welcome}
\end{itemize}
\end{frame}

\begin{frame}[fragile]{Creating Debug Ready Code}
\begin{lstlisting}[language={}, basicstyle=\ttfamily\color{codeVariable}]
g++ -g greeting.cpp -o welcome
\end{lstlisting}

\begin{itemize}
\item The \lstinline[language={}, basicstyle=\ttfamily\color{codeVariable}]{-g} flag tells the compiler to use debug info.
\item The compiled file size is much larger.
\item We may still remove this debug information using the strip command.\\
\lstinline[language={}, basicstyle=\ttfamily\color{codeVariable}]{strip welcome}
\end{itemize}
\end{frame}

\begin{frame}[fragile]{Sanitizers for Better Runtime Error Detection}

The following arguments add additional run-time checking.
\begin{lstlisting}[language={}, basicstyle=\ttfamily\small\color{codeVariable}]
g++ -g -fsanitize=return,undefined,address greeting.cpp -o welcome
\end{lstlisting}\vspace{-0.8em}
\begin{description}[\texttt{-fsanitize=undefined~~}]
  \item[\colorbox{codeBackground}{-fsanitize=return}]
     Shows error when returning without a value from non-void function.
  \item[\colorbox{codeBackground}{-fsanitize=undefined}]
    Detects some undefined behaviors.
  \item[\colorbox{codeBackground}{-fsanitize=address}]
    Detects memory addressability issues.
\end{description}\vspace{-0.8em}
  \small{See all options at \url{https://gcc.gnu.org/onlinedocs/gcc/Instrumentation-Options.html}}
\end{frame}

\begin{frame}{Adding Optimizations}
\begin{itemize}
\item The compiler can help improve the performance of your code via optimizations.
\item \lstinline[language={}, basicstyle=\ttfamily\color{codeVariable}]{g++ -O greeting.cpp -o welcome}
\item The \lstinline[language={}, basicstyle=\ttfamily\color{codeVariable}]{-O} flag tells the compiler to optimize the code.
\item Define an optimization level by adding a number to the~\lstinline[language={}, basicstyle=\ttfamily\color{codeVariable}]{-O} flag.
\end{itemize}
\end{frame}

\begin{frame}{Getting Extra Compiler Warnings}
\begin{itemize}
\item<+->
  \emph{Error messages} -- Erroneous code that does not \\ comply with the C++ standard.
\item<+->
\emph{Warnings} -- Codes that usually tend to cause errors during runtime.
\item<+->
  To receive extra compiler warnings, use the \lstinline[language={}, basicstyle=\ttfamily\color{codeVariable}]{-Wall -Wextra} arguments.\\
  \lstinline[language={}, basicstyle=\ttfamily\color{codeVariable}]{g++ -Wall -Wextra greeting.cpp -o welcome}
  \begin{itemize}
  \item
    Useful to improve the quality of our source code
  \item
    Expose bugs that will really bug us later
  \end{itemize}
\end{itemize}
\end{frame}

\begin{frame}[fragile, t]{Even More Compiler Warnings}

\begin{lstlisting}[language=Bash, deletekeywords={test}, basicstyle=\small\ttfamily\color{codeVariable}]
g++ -Wall -Wextra -Wpedantic -Wconversion -Wshadow -Wnull-dereference -Wold-style-cast greeting.cpp -o welcome
\end{lstlisting}

\begin{overlayarea}{\linewidth}{7\baselineskip}
\begin{description}[\texttt{-Wnull-dereference~}]
  \alt<1>{
    \item[\colorbox{codeBackground}{-Wall}] all the warnings about constructions that some users consider questionable, and that are easy to avoid
    \item[\colorbox{codeBackground}{-Wextra}] some extra warning flags that are not enabled by \lstinline[language={}, basicstyle=\ttfamily\color{codeVariable}]{-Wall}
    \item[\colorbox{codeBackground}{-Wpedantic}] warnings demanded by strict ISO C++
  }{
    \item[\colorbox{codeBackground}{-Wconversion}] implicit casts that may alter a value
    \item[\colorbox{codeBackground}{-Wshadow}] when a local variable/type declaration shadows another variable, parameter, type, class member, etc.
    \item[\colorbox{codeBackground}{-Wnull-dereference}] \textit{some} null pointer dereferencing
    \item[\colorbox{codeBackground}{-Wold-style-cast}] warns about use of C-style casts
  }
\end{description}
\onslide<2>{}
\end{overlayarea}

See all options at \small{\url{https://gcc.gnu.org/onlinedocs/gcc/Warning-Options.html}}

\end{frame}

\begin{frame}{Compiling Multi-Source Programs}
\begin{itemize}
\item
  Compile them with:\\
  \lstinline[language={}, basicstyle=\ttfamily\color{codeVariable}]{g++ main.cpp a.cpp b.cpp -o hello}
\item Comments:
  \begin{itemize}
  \item If external symbols are needed, use the \lstinline{extern} keyword.
  \item Source file order becomes important.
  \item As program size increases, so does compilation time.
  \end{itemize}
\end{itemize}
\end{frame}

\begin{frame}{Processing a C++ Program}
\center
\begin{tikzpicture}[thick, every node/.style={node distance=4mm, minimum height=7mm},
  box/.style={draw, fill=CSUcyan, text=black, minimum width=2.9cm}
  ]
  % Start block
  \node[box] (block1) {Source Code};
  \node[left=of block1] {Step 1};

  \node[box, below=of block1] (block2) {Preprocess};
  \node[left=of block2] {Step 2};

  \node[box, below=of block2] (block3) {Compile};
  \node[left=of block3] {Step 3};

  \node[box, below=of block3] (block4) {Link};
  \node[left=of block4] {Step 4};

  \node[box, right=7mm of block4] (blockLibs) {Libraries};

  \node[box, below=8mm of block4] (block5) {Load};
  \node[left=of block5] {Step 5};

  \node[box, below=of block5] (block6) {Execute};
  \node[left=of block6] {Step 6};

  % Arrows
  \draw[-stealth, very thick] (block1) edge (block2);
  \draw[-stealth, very thick] (block2) edge (block3);
  \draw[-stealth, very thick] (block3) edge (block4);
  \draw[-stealth, very thick, dashed] (block4) edge (block5);
  \draw[-stealth, very thick] (blockLibs) edge (block4);
  \draw[-stealth, very thick] (block5) edge (block6);

  \draw [stealth-, very thick] (block1) -- ++ (5cm, 0) |- node[anchor=south east]{Syntax Error} (block3);

  \draw [stealth-, very thick] (block1) -- ++ (5.5cm,0) |- node[anchor=south east]{Semantic Error} (block6);
\end{tikzpicture}
\end{frame}

\begin{frame}[fragile]{Selective Recompilation After Isolated Changes}
\begin{itemize}
\item With the previous command, all source files are always recompiled, even when only one has changed.
\item To overcome, we compile in multiple steps.
\begin{lstlisting}[language=bash, basicstyle=\ttfamily\color{codeVariable}]
g++ -c main.cpp # Create main.o
g++ -c a.cpp    # Create a.o
g++ -c b.cpp    # Create b.o
g++ main.o a.o b.o -o hello # link files
\end{lstlisting}

\item \lstinline[basicstyle=\ttfamily\color{codeVariable}]{-c} tells the compiler to only create an object file.
\item The final command links the objects into an executable.
\end{itemize}
\end{frame}

\section{Makefile}


\begin{frame}{Automating Program Compilation}
\begin{itemize}
\item A \emph{makefile} is a collection of instructions that \\ should be used to compile your program.
\item Once you modify some source files and type the command \texttt{make} (or \texttt{gmake} if using GNU's make), your program will be recompiled using as few compilation commands as possible.
\end{itemize}
\end{frame}


\begin{frame}[fragile]{Makefile Structure}
\begin{itemize}
\item Variable Definitions -- Define values for variables for reuse.
\begin{lstlisting}[language=make]
CPPFLAGS  = -Wall -Wextra -Wconversion -Wshadow -Wpedantic
CPPFLAGS += -g -fsanitize=return,undefined,address
SRCS = main.cpp file1.cpp file2.cpp
CXX = g++
\end{lstlisting}
\end{itemize}
\end{frame}

\begin{frame}[fragile]{Makefile Structure}
\begin{itemize}[<+->]
\item Dependency Rules -- Define the conditions a given file needs to be recompiled, and how to compile it.
\begin{lstlisting}[language=make]
main.o: main.cpp
    g++ ${CPPFLAGS} -c main.cpp
\end{lstlisting}

\begin{itemize}
\item Recompile if any of the files after a \lstinline{:} change.
\item You \textit{must} use hard tabs in makefiles! (Spaces will NOT work.)
\item \lstinline[language=make]{# is a comment}
\end{itemize}
\end{itemize}
\end{frame}

\lstset{language=make}

\begin{frame}[fragile]{Single-Source Makefile Example}
\begin{lstlisting}[basicstyle=\footnotesize\ttfamily, language=make]
# First,  list your variable(s)
CXX = g++
# First rule, which creates the program.
#    By convention, the first rule is usually all.
all: main
# compiling the source file, main.o depends on main.c
main.o: main.cpp
        ${CXX} -g -Wall -Wextra -c main.cpp
# ${CXX} uses the value of CXX variable, case sensitive
# linking the program, the program name is main
main: main.o
        ${CXX} -g main.o -o main
# cleaning everything that can be recreated with "make"
# (basically, objects, the executable, and temp files).
clean:
        rm -f main main.o
\end{lstlisting}
\end{frame}

\begin{frame}[fragile]{Multi-Source-File Example}
\begin{lstlisting}[basicstyle=\footnotesize\ttfamily, language=make]
CPPFLAGS = -g -Wall -Wextra -Wconversion -Wshadow -Wpedantic
# Top-level rule to compile the whole program.
all: prog
# The program is made of several source files.
prog: main.o file1.o file2.o
    g++ main.o file1.o file2.o -o prog
# Rule for file "main.o".
main.o: main.c file1.h file2.h
    g++ ${CPPFLAGS} -c main.cpp
# Rule for file "file1.o".
file1.o: file1.c file1.h
    g++ ${CPPFLAGS} -c file1.cpp
# Rule for file "file2.o".
file2.o: file2.cpp file2.h
    g++ ${CPPFLAGS} -c file2.cpp
# Rule for cleaning files generated during compilations.
clean:
    rm -f prog main.o file1.o file2.o
\end{lstlisting}
\end{frame}

\begin{frame}[fragile]{Phony Targets}

\begin{itemize}
\item
  Make assumes that all the targets are files.
\item
  We may want to run commands that do not represent files.\\
\begin{lstlisting}[basicstyle=\footnotesize\ttfamily, language=make]
# Rule for cleaning files generated during compilations.
clean:
    rm -f prog main.o file1.o file2.o
\end{lstlisting}

\item
  We can identify special targets that are not files.
\begin{lstlisting}[basicstyle=\footnotesize\ttfamily, language=make]
# Allow the clean target to run even if there is a file named "clean"
.PHONY: clean
\end{lstlisting}
\end{itemize}
\end{frame}

\begin{frame}[fragile]{How can we put the object files in a subdirectory?}

\pause
\begin{lstlisting}[basicstyle=\footnotesize\ttfamily, language=make, firstnumber=6]
# Make an object folder
obj:
  mkdir -p obj

# The thing to the right of | is an order-only rule
obj/main.o: main.cpp | obj
  g++ ${CPPFLAGS} -c main.cpp -o obj/main.o

obj/a.o : a.cpp | obj
  g++ ${CPPFLAGS} -c a.cpp -o obj/a.o

obj/b.o : b.cpp | obj
  g++ ${CPPFLAGS} -c b.cpp -o obj/b.o

main: obj/main.o obj/a.o obj/b.o
  g++ ${CPPFLAGS} obj/main.o obj/a.o obj/b.o -o main
\end{lstlisting}
\end{frame}

\begin{frame}{Multi-Source Make}
\begin{itemize}
\item Commands can be anything; usually they are \\\texttt{g++} commands to compile or link.
\item Commands can be multiline, use tabs
\item Other tools:
\begin{description}[GNU Autoconf]
\item [{makedepend}] Finds dependencies for your program.
\item [{GNU Autoconf}] Generates a script for your project (\lstinline{./configure}) that detects libraries your program may use and generates a make file.
\end{description}
\item We are going to focus only on \texttt{make} for this class.
%\item make param:
%\begin{itemize}
%\item Runs the dependency for param
%\item No param runs the first label found (usually all)
%\item make clean deletes files created by make
%\end{itemize}
\end{itemize}
\end{frame}

%\begin{frame}[fragile]{makedepend}
%\begin{itemize}
%\item Program to generate dependencies
%\item Parameters are C/C++ files to scan
%\item It processes \# directives to check includes
%\item Dependencies are critical, can do dependency graphs
%\item Can be added to Makefile:\\
%\begin{lstlisting}
%depend: Makefile $(SRC)
        %makedepend $(INCLUDES) $(SRC)
%## end of Makefile
%# DO NOT DELETE THIS LINE -- make depend depends on it.
%
%\end{lstlisting}
%\item Now make depend will run makedepend and add its results directly to Makefile.  It doesn't generate the commands, but it does take care of finding dependencies for you.
%\item Many compilers automatically run dependency checks every time-never misses a dependency, but wasted work
%\end{itemize}
%\end{frame}

%\begin{frame}{configure}
%\begin{itemize}
%\item Configure is part of the autoconf GNU system
%\item It makes Makefiles (ha!)
%\item Finds things (like say tcl) and updates Makefile variables with their locations
%\item Good for multi-platform Makefiles
%\item Not needed for this class
%\end{itemize}
%\end{frame}


% \section{Debugging}

% \begin{frame}{Debugging}
% \begin{itemize}
% \item Overall being about to see the entire program and use print statements is the best debugger tool.
% \item Nonetheless, there will be times where you simply can't see the whole picture...
% \item A debugger is an aid, it will help you see the whole picture.
% \end{itemize}
% \end{frame}

% \begin{frame}[fragile]{An example}
% \begin{lstlisting}
% #include <iostream>

% int get_sum(int ary[1000]) {
%     int sum = 0;
%     for (int i = 0; i <= 1000; i++) {
%         sum += ary[i];
%     }
%     return sum;
% }

% int main(int argc, char *argv[]) {
%     int ary[1000];
%     for (int i = 0; i < 1000; i++) {
%         ary[i] = i;
%     }
%     std::cout << get_sum(ary) << "\n";
% }
%\end{lstlisting}
%See the issue?  If yes/no lets use the debugger to help.
%\end{frame}
%
%\begin{frame}{Debugging C/C++ Prrograms}
%\begin{itemize}
%\item Before invoking the debugger, make sure you compiled your program with the "-g" flag (for gcc or g++) \\
%g++ -g debug\_me.cpp -o debug\_me \\
%nemiver debug\_me
%\end{itemize}
%\includegraphics[width=1.0\textwidth]{nemiver1.png}
%\end{frame}

%\begin{frame}{Let go through an example}
%\begin{itemize}
%\item nemiver demonstration.
%\end{itemize}
%\end{frame}

\section{System Architecture}


\begin{frame}{System Calls And Device Drivers}
\begin{itemize}
\item System calls access and control files and devices.
\item At the heart of the operating system (the \emph{kernel})\\
  are some device drivers.
\item The low-level functions used to access the device drivers, the system calls, include:
\begin{description}[close~]
\item[{open}] Open a file or device
\item[{read}] Read from an open file or device
\item[{write}] Write to a file or device
\item[{close}] Close the file or device
\item[{ioctl}] Pass control information to a device driver
\end{description}
\end{itemize}
\end{frame}

\begin{frame}{System Calls And Device Drivers}
  Using low-level system calls directly for input and \\output can be very inefficient.\vspace{-0.6em}
\begin{columns}
\column{0.4\textwidth}
  \begin{itemize}
  \item Why?
  \begin{itemize}
  \item Performance penalty in making a system call.
  \item The hardware has limitations
  \end{itemize}
  \item Standard libraries provide a higher-level interface to devices and disk files.
  \end{itemize}
\column{0.5\textwidth}\vspace{1em}
  \begin{tikzpicture}[thick, every node/.style={node distance=4mm, minimum height=7mm}]
    % Start block
    \node[,
      minimum width=4.4cm] (user) {User Program};
    \node[text=black, fill=CSUcyan,
      minimum width=3cm] (library) [below=of user.west, anchor=north west] { Library };

    \begin{scope}[on background layer, thick]
      \node [draw, fit={(user) (library)}, inner sep=0pt, fill=CSUslate, thick] (userAll) {};
    \end{scope}
    \node[right=of userAll] (userSpace) {User Space};


    \node[text=black, fill=CSUcyan,
      below=10mm of userAll,
      minimum width=4.4cm] (sys) {System Calls};

    \node[
      below=of sys.center,
      minimum width=4.4cm] (kernel) {Kernel};

    \node[text=black, fill=CSUcyan,
      below=of kernel.west, anchor=north west,
      minimum width=3cm] (dev) {Device Drivers};

    \begin{scope}[on background layer]
      \node[draw, fit={(sys) (dev) (kernel)}, inner sep=0pt, fill=CSUslate, thick] (systemAll) {};
    \end{scope}
    \node[right=of systemAll] (kernelSpace) {Kernel Space};


    \node[draw,
      below=of systemAll,
      minimum width=4.4cm, align=center, fill=CSUslate] (hardware) {Hardware Devices};

    % Arrows
    \draw[-stealth, very thick, shorten >=1pt] (userAll) edge node[anchor=west]{Calls} (sys);
    \draw[-stealth, very thick, shorten >=1pt] (systemAll) edge (hardware);
    \draw[-stealth, very thick, shorten >=1pt] (kernelSpace) edge (systemAll);
    \draw[-stealth, very thick, shorten >=1pt] (userSpace) edge (userAll);

  \end{tikzpicture}
\end{columns}
\end{frame}


\section{Valgrind}

\begin{frame}{What is Valgrind?}
\begin{itemize}
\item A tool to perform:
\begin{itemize}
\item memory debugging,
\item memory leak detection, and
\item memory profiling.
\end{itemize}
\item Valgrind accomplishes this by running your program inside of its virtual machine and capturing all your memory accesses/requests.
\end{itemize}
\end{frame}

\begin{frame}[fragile]{Memory Debugging: A Successful Run}
\begin{lstlisting}[basicstyle=\tiny\ttfamily\bfseries,identifierstyle=\color{white}]
==2231== Memcheck, a memory error detector
==2231== Copyright (C) 2002-2013, and GNU GPL'd, by Julian Seward et al.
==2231== Using Valgrind-3.10.0 and LibVEX; rerun with -h for copyright info
==2231== Command: ./a.out
==2231==
==2231==
==2231== HEAP SUMMARY:
==2231==     in use at exit: 0 bytes in 0 blocks
==2231==   total heap usage: 0 allocs, 0 frees, 0 bytes allocated
==2231==
==2231== All heap blocks were freed -- no leaks are possible
==2231==
==2231== For counts of detected and suppressed errors, rerun with: -v
==2231== ERROR SUMMARY: 0 errors from 0 contexts (suppressed: 0 from 0)
\end{lstlisting}
\begin{itemize}\vspace{-1em}
\item Notice Valgrind keeps track of:
\begin{itemize}
\item (heap) memory used on exit
\item How much heap memory was allocated \& deallocated.
\item How many memory errors (out-of-bounds memory access) were detected.
\end{itemize}
\item \lstinline{2231} is the process ID, which is unimportant for this course.
\end{itemize}
\end{frame}

\begin{frame}[fragile]{Memory Leak Detection}
Consider this obvious memory leak:

\begin{columns}
\column{0.46\textwidth}
\begin{lstlisting}[basicstyle=\footnotesize\ttfamily,language=C++]
int main()
{
  int *pArr = new int[512];

  return 0;
}
\end{lstlisting}

\column{0.54\textwidth}
\lstinline[basicstyle=\small\ttfamily]{g++ -g test.cpp && valgrind ./a.out}
\begin{lstlisting}[basicstyle=\tiny\ttfamily\bfseries,identifierstyle=\color{white}]
==2476== HEAP SUMMARY:
==2476==     in use at exit: 4,096 bytes in 1 blocks
==2476==   total heap usage: 1 allocs, 0 frees, 4,096 bytes allocated
==2476==
==2476== LEAK SUMMARY:
==2476==    definitely lost: 4,096 bytes in 1 blocks
==2476==    indirectly lost: 0 bytes in 0 blocks
==2476==      possibly lost: 0 bytes in 0 blocks
==2476==    still reachable: 0 bytes in 0 blocks
==2476==         suppressed: 0 bytes in 0 blocks
\end{lstlisting}
\end{columns}

\begin{itemize}
\item \emph{Definitely lost} means \textellipsis{} we lost \textellipsis{}
\item \emph{Indirectly lost} means we lost it, and it could be hard to find.
\item \emph{Possibly lost} means Valgrind was not able to determine if the memory was deallocated or not.
\item \emph{Still reachable} means you have a dangling pointer.
\end{itemize}
\end{frame}

\begin{frame}[fragile]{Memory Usage Checking}
A better version:

\begin{columns}
\column{0.46\textwidth}
\begin{lstlisting}[basicstyle=\footnotesize\ttfamily,language=C++]
int main()
{
  int *ints = new int[1024];

  delete[] ints;
  return 0;
}
\end{lstlisting}
\column{0.54\textwidth}
\lstinline[basicstyle=\small\ttfamily]{g++ -g test.cpp && valgrind ./a.out}
\begin{lstlisting}[basicstyle=\tiny\ttfamily\bfseries,identifierstyle=\color{white}]
==2591== HEAP SUMMARY:
==2591==     in use at exit: 0 bytes in 0 blocks
==2591==   total heap usage: 1 allocs, 1 frees, 4,096 bytes allocated
==2591==
==2591== All heap blocks were freed -- no leaks are possible
\end{lstlisting}
\end{columns}


Notice that Valgrind can tell us how much heap memory we are using even though we freed (deallocated) it.
\end{frame}

\begin{frame}[fragile]{Memory Access Error Detection}
Now let's use Valgrind to detect an off-by-one memory error.

\begin{columns}
\column{0.5\textwidth}
\begin{lstlisting}[basicstyle=\footnotesize\ttfamily,language=C++]
int main() {
  int *ints = new int[1024];

  for (int i = 0; i <= 1024; i++) {
    ints[i] = i;
  }

  delete[] ints;
  return 0;
}
\end{lstlisting}

Do you see the error?

\pause
\column{0.50\textwidth}
\begin{lstlisting}[basicstyle=\tiny\ttfamily\bfseries,identifierstyle=\color{white}]
==2615== Invalid write of size 4
==2615==    at 0x400673: main (in ~/t/a.out)
==2615==  Address 0x5a03040 is 0 bytes after a block of size 4,096 alloc'd
==2615==    at 0x4C298A0: operator new[](unsigned long) (vg_replace_malloc.c:389)
==2615==    by 0x40064E: main (in ~/t/a.out)
==2615==
==2615==
==2615== HEAP SUMMARY:
==2615==     in use at exit: 0 bytes in 0 blocks
==2615==   total heap usage: 1 allocs, 1 frees, 4,096 bytes allocated
\end{lstlisting}

Well, that tells us the error, but where is it?

\end{columns}
\end{frame}

\begin{frame}[fragile]{Getting A Little More Help}

Compile with debug flags! \colorbox{black}{\texttt{\color{codeVariable}g++ \textcolor{white}{\textbf{\textit{-g}}} test.c \&\& valgrind ./a.out}}

\begin{lstlisting}[basicstyle=\footnotesize\ttfamily,identifierstyle=\color{white}]
==2634== Invalid write of size 4
==2634==    at 0x400673: main (test.c:7)
==2634==  Address 0x5a03040 is 0 bytes after a block of size 4,096 alloc'd
==2634==    at 0x4C298A0: operator new[](unsigned long) (vg_replace_malloc.c:389)
==2634==    by 0x40064E: main (test.c:4)
==2634==
==2634==
==2634== HEAP SUMMARY:
==2634==     in use at exit: 0 bytes in 0 blocks
==2634==   total heap usage: 1 allocs, 1 frees, 4,096 bytes allocated
==2634==
==2634== All heap blocks were freed -- no leaks are possible
\end{lstlisting}
\end{frame}

\begin{frame}{Conclusion}
\begin{itemize}
\item Valgrind can be used to detect memory leaks and\\ memory access errors.
\item Valgrind provides other tools to profile memory usage and cache usage, but those are beyond the scope of this class.
\item Why doesn't C++ does not have built-in memory debugging like Java?
\end{itemize}
\end{frame}

\section{In the Future}

\begin{frame}{Gnuplot: Generates Simple Graphs}
\begin{center}\large
Gnuplot command-line graphing utility for Linux.

We will use it soon to observe the performance of our data structures and algorithms.
\end{center}
\end{frame}

% \begin{frame}{Profiling}
% \begin{itemize}
% \item Profiling tells you where your program is spending the most time.
% \item With that knowledge you can focus your energy on certain parts to speed up performance.  (Why?)
% \end{itemize}
% \begin{itemize}
% \item perf (hopefully)
% \end{itemize}
% \end{frame}

\section*{Lab 02}

\begin{frame}{Git: Adding an Upstream Repository}
\center{}
  Is your forked repository missing the updates\\from the course repository?\\[1em]

  \begin{tikzpicture}[thick, every node/.style={node distance=4mm, minimum height=7mm}]
    % Start block
    \node[text=black, fill=CSUcyan,
      minimum width=3cm] (upstream) {upstream};

    \node[text=black, fill=CSUcyan,
      minimum width=3cm] (origin) [right=20mm of upstream.east, anchor= west] { origin };

    \node [fit={(upstream) (origin)}, inner sep=1pt] (userAll) {};

    \node[text=CSUgold!60!white] (githubText) at ($(upstream.north)!0.5!(origin.north)$)[yshift=5pt, anchor=south] { \textbf{GitHub} };

    \begin{scope}[on background layer]
    \node [draw=CSUgold, thick, fit={(userAll) (githubText)}, fill=CSUslate, inner sep=1pt] (github) {};
    \end{scope}

    \node[text=black, fill=CSUcyan,
      minimum width=3cm, draw=CSUgold, thick] (local) [below=20mm of origin.south, anchor=north] { local };

    % Arrows
    \draw[-stealth, shorten >=1pt, very thick, color=white] (upstream) edge node[text=white, anchor=south]{fork} (origin);
    \draw[-stealth, shorten >=1pt, very thick] (origin)  to[bend right] node[anchor=east,align=right]{clone/pull} (local);
    \draw[-stealth, shorten >=1pt, very thick] (local)  to[bend right] node[anchor=west]{push} (origin);

    \onslide<2->{\draw[-stealth, shorten >=1pt, very thick] (upstream)  to[bend right] node[anchor=east]{pull} (local);}

  \end{tikzpicture}

\end{frame}

\begin{frame}[fragile]{Git: Adding an Upstream Repository}
%Add the course repository as a new remote to your local copy.

\begin{lstlisting}[language=bash, basicstyle=\small\ttfamily,breakindent=2em, identifierstyle=\color{white}]
# See your current remote repositories
git remote -v

# Add the course repository as a remote repository
git remote add upstream https://github.com/csu-cs/CSCI-315-2026-Fall.git

# Set the merging method for divergent branches
git config pull.rebase false # merge

# Get the latest updates for the class repository
git pull upstream master

# Save these changes on your GitHub fork
git push
\end{lstlisting}

\end{frame}


\begin{frame}{Lab 2: Makefiles}
\begin{center}\Large
Let's look at the lab based on today's lecture.
\end{center}
\end{frame}

\end{document}
